\problem{2.22}{Three delta-function integrals}
The sifting identity is
$\int_I f(x)\delta(x-x_0)\dd x=f(x_0)$
if $x_0$ lies in the interior of $I$, and zero if it
lies outside. Consequently
\begin{align}
\int_{-3}^{1}(x^3-3x^2+2x-1)\delta(x+2)\dd x
&=(-2)^3-3(-2)^2+2(-2)-1=\boxed{-25},\\
\int_0^\infty[\cos(3x)+2]\delta(x-\pi)\dd x
&=\cos(3\pi)+2=\boxed{1},\\
\int_{-1}^{1}\ee^{|x|+3}\delta(x-2)\dd x
&=\boxed{0}.
\end{align}
The third integral vanishes because the delta's support
lies outside the domain, regardless of the integrand
evaluated at its center.

\problem{2.23}{Scaling of the delta and distributional derivative of a step}
For nonzero real $c$ and a smooth test function $f$,
substitute $u=cx$, taking account of orientation if
$c<0$:
\begin{equation}
\int_{-\infty}^{\infty}f(x)\delta(cx)\dd x
=\frac1{|c|}\int f(u/c)\delta(u)\dd u
=\frac{f(0)}{|c|}.
\end{equation}
Therefore $\boxed{\delta(cx)=\delta(x)/|c|}$ as
distributions. For the Heaviside step $\theta$,
integration by parts against a compactly supported
smooth test function gives
\begin{equation}
\langle\theta',f\rangle
=-\int_0^\infty f'(x)\dd x=f(0)
=\langle\delta,f\rangle,
\end{equation}
and hence $\boxed{\theta'(x)=\delta(x)}$
in the distributional sense. The conventional
value $\theta(0)=1/2$ has no effect on this result.

\problem{2.24}{Uncertainty product for a delta-bound state}
For $V=-\alpha\delta(x)$ with $\alpha>0$, define
$\kappa=m\alpha/\hbar^2$. The normalized bound state is
\begin{equation}
\psi(x)=\sqrt\kappa\,\ee^{-\kappa|x|},
\qquad E=-\frac{\hbar^2\kappa^2}{2m}.
\end{equation}
Symmetry gives $\avg x=\avg p=0$ and
\begin{equation}
\avg{x^2}=2\kappa\int_0^\infty
x^2\ee^{-2\kappa x}\dd x=\frac1{2\kappa^2}.
\end{equation}
The first derivative is
$\psi'=-\kappa\,\operatorname{sgn}(x)\psi$,
whereas the second derivative contains a delta:
\begin{equation}
\psi''=\kappa^2\psi-2\kappa\psi(0)\delta(x).
\end{equation}
Therefore, using $\avg{p^2}=-\hbar^2\int\psi^*\psi''\dd x$,
\begin{equation}
\avg{p^2}=-\hbar^2\kappa^2+
2\hbar^2\kappa|\psi(0)|^2
=\hbar^2\kappa^2.
\end{equation}
Consequently
\begin{equation}
\boxed{\sigma_x=\frac1{\sqrt2\kappa},\quad
\sigma_p=\hbar\kappa,\quad
\sigma_x\sigma_p=\frac{\hbar}{\sqrt2}>
\frac{\hbar}{2}}.
\end{equation}
\check{The delta term is essential; ignoring the derivative
jump would incorrectly produce a negative $\avg{p^2}$.}

\problem{2.25}{Bound-scattering orthogonality in the delta well}
Let $\psi_b$ be the normalized delta-well bound state
with $E_b<0$ and let $\psi_k$ be a generalized scattering
eigenfunction with $E_k>0$. Both satisfy the same matching
condition at $x=0$:
$\psi'(0^+)-\psi'(0^-)=-2m\alpha\psi(0)/\hbar^2$.
Multiply the equations $\hat H\psi_b=E_b\psi_b$
and $\hat H\psi_k=E_k\psi_k$ by $\psi_k^*$
and $\psi_b$, subtract and integrate on the two
half-lines. The Wronskian boundary contributions at
$0$ cancel because both states obey the same jump,
and the terms at infinity vanish because
$\psi_b,\psi_b'\to0$ exponentially and $\psi_k$ is bounded.
Thus
\begin{equation}
(E_k-E_b)\int_{-\infty}^\infty\psi_b(x)\psi_k^*(x)\dd x=0.
\end{equation}
Since $E_k\ne E_b$,
\begin{equation}
\boxed{\int\psi_b(x)\psi_k^*(x)\dd x=0}.
\end{equation}
\check{The scattering state is not individually in
$L^2(\mathbb R)$, but this mixed overlap is convergent.}

\problem{2.26}{Fourier transform of a delta distribution}
For the symmetric convention
$\widehat f(k)=(2\pi)^{-1/2}
\int f(x)\ee^{-\ii kx}\dd x$,
\begin{equation}
\boxed{\widehat\delta(k)=\frac1{\sqrt{2\pi}}}.
\end{equation}
Formally applying the inverse transform gives
\begin{equation}
\boxed{\delta(x)=\frac1{2\pi}
\int_{-\infty}^\infty\ee^{\ii kx}\dd k}.
\end{equation}
This is an identity of distributions, not an
ordinary convergent improper integral.
For instance, a Gaussian regulator gives
\begin{align}
\frac1{2\pi}\int\ee^{-\epsilon k^2}
\ee^{\ii kx}\dd k
&=\frac1{\sqrt{4\pi\epsilon}}
\ee^{-x^2/(4\epsilon)}
\longrightarrow\delta(x)
\quad(\epsilon\downarrow0).
\end{align}
Against any smooth, rapidly decaying test function,
the right-hand side converges to its value at $x=0$.

\problem{2.27}{Bound states of two attractive delta centers}
For $V=-\alpha[\delta(x-a)+\delta(x+a)]$
with $\kappa_0=m\alpha/\hbar^2$ and $E<0$,
set $\kappa=\sqrt{-2mE}/\hbar>0$. By parity, a
normalized bound state has the form
\begin{align}
\psi_e(x)&=\begin{cases}
C\cosh(\kappa x),&|x|<a,\\
C\cosh(\kappa a)\ee^{-\kappa(|x|-a)},&|x|\ge a,
\end{cases}\\
\psi_o(x)&=\begin{cases}
C\sinh(\kappa x),&|x|<a,\\
C\,\operatorname{sgn}(x)\sinh(\kappa a)
\ee^{-\kappa(|x|-a)},&|x|\ge a.
\end{cases}
\end{align}
Continuity at $a$ is automatic. Imposing the derivative
jump $\psi'(a^+)-\psi'(a^-)=-2\kappa_0\psi(a)$ gives
\begin{equation}
\boxed{\kappa_e=\kappa_0(1+\ee^{-2\kappa_e a}),\qquad
\kappa_o=\kappa_0(1-\ee^{-2\kappa_o a})}.
\end{equation}
The even state always exists; a positive odd root exists
only for $2\kappa_0a>1$. For the specified strengths:
\begin{center}
\begin{tabular}{lcc}
\toprule
$\alpha$&$\kappa_e a$&$\kappa_o a$\\
\midrule
$\hbar^2/(ma)$&$1.108858$&$0.796812$\\
$\hbar^2/(4ma)$&$0.369418$&none\\
\bottomrule
\end{tabular}
\end{center}
Each energy is $E=-\hbar^2(\kappa a)^2/(2ma^2)$.
As $a\to0$, the even state approaches a single
delta of strength $2\alpha$, with
$E_e\to-2m\alpha^2/\hbar^2$; the odd state disappears.
As $a\to\infty$, the two bound levels converge
to the isolated-well energy $-m\alpha^2/(2\hbar^2)$.

\problem{2.28}{Transmission through the double-delta potential}
For $k=\sqrt{2mE}/\hbar>0$ and
$\beta=m\alpha/(\hbar^2k)$, write a local plane-wave
solution as $A\ee^{\ii kx}+B\ee^{-\ii kx}$ to
the left and $F\ee^{\ii kx}+G\ee^{-\ii kx}$
to the right. Matching at an attractive delta at
$x=x_0$ yields
\begin{equation}
\begin{pmatrix}F\\G\end{pmatrix}=
\underbrace{\begin{pmatrix}
1+\ii\beta&\ii\beta\ee^{-2\ii kx_0}\\
-\ii\beta\ee^{2\ii kx_0}&1-\ii\beta
\end{pmatrix}}_{M(x_0)}
\begin{pmatrix}A\\B\end{pmatrix}.
\end{equation}
For centers at $-a$ and $a$, $M=M(a)M(-a)$.
Because $\det M=1$, left incidence ($G=0$) gives
$F/A=1/M_{22}$, where
\begin{equation}
M_{22}=(1-\ii\beta)^2+\beta^2\ee^{4\ii ka}.
\end{equation}
Thus
\begin{equation}
\boxed{T(k)=
\frac1{\left|(1-\ii\beta)^2+
\beta^2\ee^{4\ii ka}\right|^2}}.
\end{equation}
This is real, nonnegative, and no larger than one
for a real potential. As $a\to0$,
$T\to[1+4\beta^2]^{-1}$, precisely the result for
a single attractive delta of strength $2\alpha$.
